\chapter{Gender issues in translation studies} \label{chapter2-1}

This chapter outlines the theoretical background of the monograph within TS, situating the work in the field of feminist and queer translation. Section \ref{TSrelevance} discusses the relevance of the present work for the disciplines and briefly conceptualises translation as a queer and social practice. The key concepts introduced here are further elaborated on in Section \ref{translationactivism}: A brief overview of the developments within TS that led to the rise of feminist translation is traced and activism in translation is discussed. In addition, key questions from the field of gender studies, such as sex, gender and gender identity, are also explored. Attention is paid to how categories of sex and gender are socially constructed and still often conceived as binary, even though research in different fields (e.g. \cite{fausto2012sex,richards2016non}) demonstrates that there are more than two sexes and genders. An overview of the nature and types of non-binary genders is therefore provided. In Section \ref{sec:powerlanguage}, the interplay between language, society, and gender is introduced. Languages are the result of cultural conventions, and differences in their structures can lead to social and cultural differences (see e.g. \cite{prewitt2012gendering,santacreu2013female}) that both reflect and impact assumptions on sex and gender (see \cite{boroditsky2003sex,stahlberg2007representation}). Thus, it is posited that language (re)produces norms, social expectations, and power relations and can be used both to discriminate as well as to promote equality \citep{sprachhandeln2015tun}. 

\section{The social responsibility of translation} \label{TSrelevance}
In his seminal work, \textcite{venuti1995translator} uses the term “invisibility” to describe the situation and activity of translators within Anglo-American culture. Venuti illustrates how translated texts are judged by their fluency, implying that a translation is considered good when it is not perceived as such, but rather as an original work. This perspective carries significant implications. The translation is seen as a derivative, fake, and potentially false copy. Consequently, there has been a longstanding emphasis within TS on the notion of \textit{fidelity} to the source text and its author (see \cite{pym2017exploring}). Furthermore, \textcite[1]{simon1996gender} points out that the language used to describe translation often ``dips liberally into the vocabulary of sexism, drawing on images of dominance and inferiority, fidelity and libertinage''. For example, the phrase ``les belles infidèles'', once commonly used to describe translations, reflects the expectation that a target text should either be beautiful or faithful, echoing societal perceptions of women \citep[10]{simon1996gender}.

While calling into question the position of discursive inferiority of women and translations, feminist TS understand language as having a performative nature \citep[2]{simon1996gender}. The idea of performativity is also at the core of Butler's (\cite*{butler1990gender}) theorisation of gender as socially and discursively constructed. By asking how social, sexual and historical differences are expressed and how they can be transferred across languages, \textcite{simon1996gender} supports sociolinguistic views of language as actively creating meaning and contributing to reality \citep{fairclough2015language}. Being a socially and culturally embedded practice \citep{wolf2007sociology}, translation can either uphold or challenge the dominant regimes of knowledge and power. \textcite{tymoczko2010translation}, for example, outlines how evolving conceptualisations of translation from a mere substitution of linguistic code to an ethical, ideological, and political practice led to considering translators as agents for social change. Translators are sometimes even caught between the two aforementioned positions, representing both institutional power and people seeking empowerment. Translators therefore act ``as a kind of double agent in the process of cultural negotiation'' \citep[xix]{tymoczko2002introduction} and their double role might be considered a strength that is indicative of the complexity of contemporary cultural hybridity.

From a theoretical point of view, feminist and queer translation contributes, together with other paradigms such as postcolonial translation, to the deconstruction of rigid binary oppositions prevalent in TS, such as original/copy and author/translator. Translation thus becomes ``a site of struggle in the negotiation and production of meaning'' \citep[176]{spurlin2017queering} that shapes social and cultural identities. This perspective underscores the social responsibility of translators, defined by \textcite[122]{drugan2017translation} as ``action-oriented and dynamic, encompassing value judgements and decisions that may lead as much to resistance as to acceptance and commitment to sustain a form of social consensus''. However, social responsibility remains an understudied concept in relation to the translation of GFL. 

From an empirical perspective, feminist and queer translation aims to give visibility to the work of feminist and queer translators today and throughout history, address how the categories of sexual and gender identities are the results of social construction, and demonstrate how texts can be translated either to support or to question the status quo \citep{simon1996gender,von2016translation,burton2010inverting,epstein2017eradicalization}. For instance, feminist and queer translation practices can critically highlight hegemonic, sexist, homo- and transphobic views expressed in texts or ``invert'' \citep[57]{burton2010inverting} them.

The present monograph adopts a feminist and queer approach to translation, conceptualising it as a cultural, political, and ethical negotiation process. This understanding is rooted in the idea of translation as a \textit{queer practice} characterised by ``the study of borders, and their deconstruction and rearrangement'' \citep[172]{spurlin2017queering}. According to Spurlin, translation disrupts and reconfigures the imaginary boundaries between concepts such as man/woman and author/translator, replacing dichotomies with continuums or spectrums. In other words, translation creates crossings across linguistic borders and social categories, producing new, hybrid forms of meaning and new knowledge through these very encounters, which may result in questioning the actual borders of language and cultures \citep[173]{spurlin2017queering}. This is exactly what language professionals do when they translate non-binary GFL.

This work aligns with the paradigms of feminist and queer translations for different reasons. First, it investigates the translation and PE of texts referring to non-binary individuals, a marginalised group often overlooked in TS research. Second, it demonstrates how translation, as a feminist and queer practice, engages in a cultural, political, and ethical process of negotiation. Language professionals have the power to either erase gender diversity and make it fit societal gender binary norms or to represent it. Such a decision implicates social responsibility. Misgendering translations may lead to identity invalidation, with detrimental effects on the mental health of transgender individuals \citep{mclemore2015experiences,mclemore2018minority}. 

While English is mostly gender-neutral, German requires constant gender inflection and it has no equivalent to singular \textit{they} for reference to a hypothetical person, people of unknown gender, or non-binary individuals. Language professionals are thus required to experiment with different, even new and disruptive, GFL strategies. Their use carries ideological and political connotations in German that might not be present in the English source text. It might therefore be possible that language professionals have to explain and/or justify their linguistic choices. The translation process thus becomes visible, with professionals effectively engaging in an activist agenda by serving as spokespeople for a marginalised group. They mediate among and negotiate with numerous co-textual and contextual factors, such as clients’ specifications, personal and societal views and/or beliefs, and target audience amongst others.

Observing the translation and PE process and discussing it with language professionals contribute towards bringing translators to the forefront of the discourse around gender fairness in language and translation, showcasing their role as mediators and agents for social change. This is all the more important because this research expands the scope of activist, feminist, and queer translation beyond literary contexts to encompass fields like audiovisual and news translation, which have received less attention in the literature \citep{baer2018queering,von2020routledge}. 

MT is also relevant to (queer) TS. Language professionals are constantly confronted with technology and increasingly receive PE assignments. MT systems, trained on large amounts of data that reflect a majority language use, exhibit gender bias \citep{savoldi2021gender}. As a widely used translation technology, MT propagates hegemonic views on gender, discriminating against marginalised groups. It is therefore important that feminist and queer translation addresses this issue and engages in multidisciplinary endeavours to mitigate it. Furthermore, the propagation of bias in MT substantiates the idea that translators can act as agents for social change in a process of social, ethical, and political negotiation.

%----------------------------------------------------------------------
%	SECTION 1
%----------------------------------------------------------------------

\section{Shifting translation concepts: The emergence of culture and ideology} \label{translationactivism}

Gender issues have been of interest since second-wave feminism not only in the field of linguistics but also in TS \citep{simon1996gender,von2016translation}. This is also due to the development of the discipline which has emancipated itself from linguistics and flourished in a well-established research area. In the following paragraphs, I present a brief overview of the changes in the conceptualisation of translation practice that led to the emergence of paradigms such as feminist and queer translation.

In Western Europe, translation has long been defined mainly as a simple process of cross-linguistic transfer. During World War II, translation was conceptualised as an instrument to decipher the linguistic and cultural codes of both enemies and allies and as an instrument of propaganda \citep[3]{tymoczko2010translation}. This conceptualisation is also at the core of the first MT systems \citep{hutchins2015machine} and it is specifically from the interest in automated translation (and its difference from human translation) that TS emerged as an academic discipline.

At the time, TS was strongly influenced by linguistics and the debate on translation revolved around the concept of equivalence \citep[59]{munday2022introducing}. The process of translation tended to be seen as a mere substitution of a linguistic code through another and the fundamental criterion to establish the quality of translations was their fidelity and equivalence to a source text (see \cite{pym2017exploring} for an overview of translation theories and their development). During this period, the work of scholars such as \textcite{jakobson1959linguistic}, \textcite{nida1964}, and \textcite{newmark1981} was particularly influential. Jakobson was influenced by structuralism and focused on cross-linguistic differences in grammatical and lexical forms, noting that ``[l]anguages differ essentially in what they must convey and not in what they may convey'' \citep[129]{jakobson1959linguistic}. An example of such differences is gender. For instance, Romance languages must convey gender in nouns, whereas English is mostly gender-neutral. \textcite{nida1964} borrowed concepts and terminology from semantics and pragmatics, and distinguished between formal and dynamic equivalence \citep{munday2022introducing}. The first focuses on the form and content of the message, whereas the second on the effect, thus introducing a reader-based orientation to translation theory. \textcite{newmark1981} was influenced by linguistics as well and proposed two types of translation, semantic and communicative. They can respectively be compared to Nida's formal and dynamic equivalence.

By the 1970s, two different paradigms gained prominence within TS, i.e. functionalist approaches and system theories. For instance, \textcite{reiss1971moglichkeiten} expanded the concept of equivalence to the text level and focused on the analysis of text types and genres. Furthermore, \textcite{reiss1984grundlegung} and \textcite{nord1988} advocated for a more dynamic understanding of translation processes, emphasising the importance of context and purpose in translation and paving the way for a more nuanced analysis of translation practices. A further important shift in this direction was polysystem theory \citep{even-zohar1990}, which considered literature as part of the cultural, literary, and historical system of the target language.

The 1980s witnessed the rise of Descriptive Translation Studies (DTS) as a paradigmatic shift within the field. Scholars such as \textcite{toury1995descriptive} spearheaded this movement, advocating for an empirical approach to translation research that focused on describing translation practices as they occur in real-life contexts. DTS abandoned prescriptive approaches and aimed to describe translation products and texts according to their cultural and political coordinates (\cite{toury1995descriptive}, \cite[64--70]{pym2017exploring}), paving the way for considering culture as a significant analytical category together with language in the 1990s. 

At that time, \textcite{snell1990linguistic} attested to a cultural turn within the discipline, moving the conceptualisation of translation as text to translation as culture and politics \citep[198]{munday2022introducing}. Translation was accordingly redefined as a type of intercultural transfer (\cite{bassnett1990translation}, \cite{venuti2018rethinking}, \cite*{venuti1995translator}, \cite[5]{tymoczko2010translation}). This led translation theorists to reflect on the role of translation as an instrument of power and on its ideological aspects \citep[3--4]{tymoczko2010translation}. \textcite[25]{toury1982rationale}, for example, claimed that ``translations are facts of one system only: the target system'', meaning that translation strategies, norms, functions, and purposes are determined by the culture in which they are produced. Thus, translations have political and ideological implications. 

Ever since, attention has been turned to issues of power, ideology, dominance, norms, and ethics in translation. The decision to translate (or not) entire texts as well as translation shifts would reveal historical and geopolitical patterns of target cultures as well as political constraints or initiatives underlying the choices made by translators \citep[6--7]{tymoczko2010translation}. The decision to erase references to homosexuality could be interpreted as reproducing widespread homophobic views circulating in society. Conversely, the decision to translate literary works by women that are not considered part of the traditional literary canon may be seen as a political decision to oppose a patriarchal system that excludes women from many fields like art and literature. 

This cultural turn, characterised by the dismissal of linguistic theories of translation \citep{bassnett1990translation}, led to a stronger image of translators as active agents. Translation started to be seen as an instrument of resistance and activism \citep[7]{tymoczko2010translation}. Venuti (\cite*{venuti2018rethinking}, \cite*{venuti1995translator}) called translators to action, urging them to become visible, i.e. to become involved in cultural and ideological matters. Translators were regarded as ``agents for social change'' \citep[3]{tymoczko2010translation}, using translation to actively change societies, whilst translation began to be considered as an ethical, ideological, political practice -- far more than a simple matter of linguistic transfer. This conceptualisation shift set up the basis for feminist TS, as explained in Subsection \ref{feministtranslation}. Accordingly, from the late 20th century onwards, attention has been paid to the role of activist translators and their impact on society \citep{tymoczko2002introduction,tymoczko2010translation}. 

Throughout the years, activist translators have dealt with matters such as European colonialism, imperialism, the Hispanic revolution in America \citep[227]{tymoczko2010space}, sexism (\cite{von1991feminist}, \cite*{von2016translation}), and discrimination against the LGBTQIA+ community \citep{burton2010inverting}. The last two points are of particular interest for the present work, as queer translation \citep{spurlin2014queering,baer2017introduction} can be seen as an extension of feminist translation, and it represents the positioning of the present work within the framework of TS. Translators have been coping with oppression brought about by dictatorships, colonialist regimes imposing social hierarchies, military oppression as well as with revolutions, nation-building, gender liberation, and supporting shifts in cultural values \citep[229]{tymoczko2010space}. To do so, they have adopted several different strategies. 

\textcite[229--235]{tymoczko2010space} posits that activism in translation can be achieved by a wide variety of strategies, also highlighting how these are situated in a specific culture, space and time. Struggles and protests can change rapidly, so translation must be constantly adapted to shifting cultural, historical, and political constellations. Translations and strategies are therefore valid in and specific to a certain period \citep[234--235]{tymoczko2010space}. To summarise, activism in translation can be performed through:
\begin{itemize}
\setlength\itemsep{-0.5em}
    \item Zero translation: This is the refusal to translate a text or parts of it, hence refusing to transmit cultural information, e.g. in Hawaii, translations have been refused to actively oppose colonial authority \citep{aiu2010ne}.
    \item Choice and manipulation of text types and/or genres, e.g. Monteiro Lobato inserted political messages in children's books during Getúlio Vargas' dictatorship in Brazil \citep{milton2006resistant}.
    \item Text amplification: Translators add comments and paratextual materials so as to explain their political and ideological views. Such strategies have often been adopted in feminist translation (an overview can be found in \cite{von1991feminist}).
    \item Text simplification: This is often done so that texts are fit for specific purposes.
    \item The insertion or the refusal of foreign elements into a translation.
\end{itemize}
There is no single activist strategy for all situations. Instead, the selection depends on the specific context of resistance \citep[230]{tymoczko2010space}.

%----------------------------------------------------------------------
%	SUBSECTION 
%----------------------------------------------------------------------
\subsection{Feminist translation} \label{feministtranslation}

Before 1955, gender only described a grammatical feature of language \citep[89]{haig2004inexorable}. With second-wave feminism, gender began to indicate differences in men's and women's societal roles and opportunities that are the result of societal structures \citep[5]{von2016translation}. %As discussed in Chapter \ref{Sexism in language}
The feminist movement sought to counteract sexism and achieve equality by concentrating on the analysis of gender and its manifestation in language (see for example \cite{von1991feminist,von2016translation}, \cite{pauwels2003linguistic}). The concept of patriarchy has been used in feminist scholarship to describe the system that privileges men based on alleged gender differences. Feminist scholars have tried to theorise the development and the history of patriarchy \citep{miller2017patriarchy}. The literature produced in this field is vast (see again \cite{miller2017patriarchy} for an overview) and goes beyond the purpose of the present work.

Feminists felt the urge to focus on language as a powerful tool developed by men which shapes society \citep[8]{von2016translation}. Patriarchal language was seen not only as an expression of the system that oppresses women but also as playing an active role in their oppression. Therefore, it had to be studied and reformed in order for women to reach equality \citep[3]{cameron1985feminism}. Linguistic analyses by feminists concentrated on factors such as (i) the study of sex/gender differences in the use of language; (ii) the representation of women as well as sexism in language; (iii) women's feeling of alienation in language; and (iv) the relationships between language and power (\cite[7]{cameron1985feminism}, \cite[8]{von2016translation}, \cite[73--75]{kramer2016feminist}). 

In Western Europe and North America, women started to reform language in the 1970s and 1980s actively. For example, they produced handbooks of non-sexist language, rewrote dictionaries, criticised sexist language, and introduced new literary forms to represent their experiences \citep[9]{von2016translation}. In particular, feminist writers in Quebec were very active in producing experimental work that tried to challenge patriarchal language, hence giving birth to what is called \textit{écriture féminine} (female writing) \citep[72]{von1991feminist}. These feminist writers tried to develop a new language for women, thus using new words, especially to describe women's bodies and experiences, spellings, grammar, and metaphors (\cite[72]{von1991feminist}, \cite[15]{von2016translation}). Such reflections on patriarchal language, as well as experiments to create a new idiom that would make women visible and allow them to achieve equality, also had an impact on translation both as a practice and as a discipline (see \cite{von1991feminist,von2016translation}). 

For translators, this innovative work meant they had to experiment with new approaches. Female writing largely dealt with the topic of the female body. Derogatory vocabulary was reclaimed, and new words were coined to overcome objectified patriarchal descriptions and the lack of an everyday language to articulate them. Writings contained puns, cultural references, new grammatical forms, and played with sound associations \citep[15]{von2016translation}, thus forcing translators to develop creative methods to reproduce the effect of the source texts. These included active intervention in translation after considering differences in linguistic structure and cultural situations, and the use of footnotes when strategies equivalent to the source texts were not found \citep{von2016translation}. 

Feminists also started to practise translation activism. There were two main approaches at the time: interventionist translation and the translation of female literature. The former saw translators actively shape target texts to remove or alter sexist elements of the source text \citep[24--26]{von2016translation}. This process could occur more or less openly and could be more or less politically driven: \textcite[11]{lotbiniere1991re}, for instance, claims that translation is a political activity, and as such, her textual interventions are justified. \textcite[74--76]{von1991feminist} identifies four strategies that characterise feminist translation. These are:
\begin{itemize}
\setlength\itemsep{-0.5em}
    \item Supplementing: This compensates for differences between languages or adds meaning to the text. For example, noun forms for people differ based on their sex/gender in French whilst they generally do not in English. When translating from the former into the latter, sex/gender can be specified. 
    \item Footnoting: Translators stress their presence and try to convey with footnotes the multiple meanings of the source text that would otherwise be lost.
    \item Prefacing: Translators reflect and comment on their work, often explaining their political views and/or their translation strategies.
    \item Hijacking: Translators appropriate and manipulate source texts, using a mix of the above strategies as well as a didactic approach.
\end{itemize}

The second approach consisted of selecting and translating women's work that was not part of the traditional literary canon. The aim was to make their creativity available to the general public as it was maintained that patriarchal society defines aesthetics and literary values \citep[30]{von2016translation}. For this reason, canons have mainly consisted (and in part still consist) of works produced by men. Attention was paid to issues such as censorship as well as the silencing and non-recognition of women's contributions to culture and literature \citep[2]{von2011preface}. Though these early theories of feminist translation focused exclusively on women, TS have also been influenced by queer theories that seek to overcome binary categories such as those of sex/gender \citep[3]{von2011preface}. 


%----------------------------------------------------------------------
%	SUBSECTION 
%----------------------------------------------------------------------
\subsection{Queer translation}

In order to understand queer translation, a brief explanation of the terms \textit{queer} and \textit{queer theory} is required. TS has only recently embraced queer theory and there is still often conceptual confusion within the discipline \citep[1]{baer2017introduction}. Queer can have different meanings: (i) it can be used as a synonym for homosexual or also to denote a person who does not fit normative ideas of gender and/or sexuality; (ii) it was originally a derogatory term, but nowadays it is often used as a form of linguistic reappropriation; (iii) it refers to a theoretical model that rejects binary oppositions like that of homosexual/heterosexual \citep[2]{baer2017introduction}. Queer theory questions the authenticity and universality of dominant structures of power and knowledge and emphasises their constructedness. As a consequence, queer can be used as a verb in order to describe such a critical process \citep[3]{baer2017introduction}.

From the 1970s onwards, TS focused on the issues of power, ideologies and ethics underlying the translation process while concepts such as identity and sexuality began to permeate the debate around translation (cf. Section \ref{translationactivism} and Subsection \ref{feministtranslation}). In the 1990s, LGBTQ studies also started to pay attention to other discriminated minorities as well as their relationships with the majority culture. Sexuality also began to be conceptualised as fluid \citep[2f]{amenta2020tradurre}. This led to a queer turn in TS.

This new paradigm in TS has helped bring into question classical translation theories and approaches, also because it rejects binary categories such as those of source and target texts that have shaped the conception of the discipline \citep[3]{amenta2020tradurre}. Both queer theory and TS can benefit from their intersection, as Baer \& Kaindl explain:
\begin{quote}
     To the extent that queer theory problematizes the representation of otherness, and translation studies highlights the otherness inherent in representation, bringing together queer theory and translation studies should productively destabilize not only traditional models of representation, understood as mimesis, reflection, and copying, but also the authorial voices and subjectivities they project. \citep[1]{baer2017introduction}
\end{quote}
TS can finally overcome the dichotomy between the authoritative source and the subordinate target text, which led to the general view that translation must be faithful to the original, master text. Queer theory, as the study of borders and their deconstruction and/or rearrangement, promotes a radical rethinking of translation as a complex cultural, political, and ethical process of negotiation \citep{spurlin2014queering}. At the same time, the Western-centric view of queerness in queer studies can be dismantled, showing how gender and sexuality are differently constructed across language, cultures, and history \citep{spurlin2014queering}. These concepts are explored in more detail below, but first, a definition of queer texts and an account of the translation problems they present are provided.

Queer texts are characterised by the presence of non-normative sexual identities that deconstruct gender and sexual norms. The following challenges arise when translating them \citep[5]{amenta2020tradurre}:
\begin{itemize}
\setlength\itemsep{-0.5em}
    \item Sociolects: the use of camp talk by queer characters, or of female nicknames by gay men to refer to other gay men, neologisms, and semantic variations that are typical of the language of queer characters.
    \item Unusual grammatical and syntactical structures: often used to deconstruct the gender binary. For example, masculine and feminine gender markers may be alternated in a text to refer to the same person. When gender systems differ across languages, translators have to use synonyms or find compensating solutions, such as making gender explicit in other parts of the text.
    \item Register and style: for instance, the use of double entendres and inflections in speech.
    \item Intertextuality: allusions to a character's or the author's sexuality can be made via reference to other texts, thus challenging the reader and the translator, who must recreate such allusions or substitute references to produce a similar effect. 
\end{itemize}

According to \textcite{demont2017three}, there are three different approaches to the translation of literary queer texts, namely, misrecognising, minoritising, and queer translation. The misrecognising translation ignores queer elements of the source text and suppresses its disruptiveness -- this means that queer characters are often erased from texts and translations are ``straightened''. The minoritising one reduces connotative elements to their denotative equivalents, thus making queerness explicit but flattening the power of queer texts which are often based on ambiguity, multiple meanings, and meaning associations. Finally, there is the proper queer translation -- i.e. as intended for example by \textcite{burton2010inverting}, which can consist of two different approaches. The first is the comparison and critique of a source text and its translation to investigate whether queer elements have been concealed or erased in the target texts. The second tries to find strategies to maintain the source text's queerness. 

\textcite{epstein2017eradicalization} distinguishes between the queer translation of queer and homophobic texts. In the first case, similarly to feminist translation, strategies such as supplementing, prefacing, footnoting, or radical changes can be applied to reproduce the queerness of the source text. In the second case, texts can be translated to reveal the homophobic, biphobic, or transphobic views a text contains. Such a practice in particular is proposed by \textcite[57]{burton2010inverting}, who uses the term \textit{inversion}, meaning ``a turning of the text against itself: inverting the hidden power relations of heterosexism by revealing and underscoring them through techniques borrowed from feminists''. Queer translation is therefore ideologically and methodologically related to feminist translation. 

In addition, queer translation is ``antihomophobic in motivation and practice, and destabilising and historicising of gender, sex, and sexuality norms'' \citep[55]{burton2010inverting}. Since the 1970s, discourses on discrimination and sexism have broadened and new identities have gained visibility. Debates on how translation can reinforce or counteract homophobic views started to emerge accordingly \citep[54]{burton2010inverting}. Queer translation also adopts an intersectional lens (see \ref{languapepowerdiscrimination} for a definition), i.e. it considers how queerness and other dimensions such as race and ethnicity interact with gender and sexuality \citep[55]{burton2010inverting}.

Since translation allows the rendering of texts from different cultures and ages in which queerness is differently conceptualised, it can also foster the production of new knowledge in the field of identities and sexuality. Originally, the term \textit{queer} is anglophone and as a concept it is related to Western cultures. However, queerness is differently conceived, experienced and lived in other cultures, and so it was in other ages. Therefore, it is also differently described. Many non-Western experiences of queerness are reduced to Western equivalents or completely erased in translation. Translation as a queer practice can thus help reduce ``the monolingual, anglophone biases'' of both TS and queer theory \citep[299]{spurlin2014queering}: the analysis of other gender and sexual identities in translation may produce new knowledge that stimulates shifts in social identities and categories. Such an analysis shifts the attention to the way gender and sexuality are represented through language(s). 

Analyses on how gender and sexuality differ across languages, on changes of meaning as well as on gaps between cultures can impact translation itself and how it is conceived. With regard to translation theory metaphors, \textcite[301f]{spurlin2014queering} notes how the translated text has historically been conceived as female, that is, subordinated to the authoritative original. In the field of queer translation, such feminisation of the target text is replaced by a (re)gendering that overcomes the source/target binary (see \cite{chamberlain1988gender} for metaphors used to describe translation). Since translation is an active process that involves language and culture as well as ideological and political aspects, the original is not reproduced but rather transformed. As a matter of fact, translation is a process of negotiation and meaning production, a site of struggle that questions the relationship between the original text, i.e. the dominant male, and the translated text, i.e. the submissive female \citep[206]{spurlin2014gender}. Queering translation shows the power of translation to mitigate asymmetrical relationships between languages and cultures that reproduce existing power relations \citep[213]{spurlin2014gender}. Hence, cultural and historical differences in source texts should be preserved as a way to oppose monolingualism and ethnocentrism \citep{spurlin2014gender}, often intended as a tendency to adapt source texts to Western, and especially anglophone cultures.

%----------------------------------------------------------------------
%	SUBSECTION 
%----------------------------------------------------------------------
\subsection{Translating gender: Challenges} \label{sub:translatinggender}

Subsection \ref{sub:language/gender} addresses in more detail the extent of complexity in the relationship between gender and language. Here, it is sufficient to briefly highlight that (i) linguistic elements are used to refer to the extralinguistic reality, in this case gender identity; (ii) linguistic structures to do so vary across languages; (iii) linguistic gender has also social connotations. Accordingly, TS have analysed such issues in translation, highlighting for example how social gender influences translation choices \citep{jakobson1959linguistic,nissen2002aspects,di2020grammatical}. This section presents a concise overview of some of the problems arising from translation among languages with different gender systems along with possible solutions. 

Translation from a grammatical gender language, i.e. when every noun has a gender, into a notional gender language, i.e. when pronouns only and few nouns have a gender (see Subsection \ref{sub:language/gender} for a classification of languages based on their gender system and an in-depth explanation), can be particularly difficult as their agreement structures differ greatly \citep[26]{nissen2002aspects}. For instance, when translating adjectives and nouns describing a person from Spanish into English, gender information is lost. This can be compensated by lexically marking the terms in question by adding a word specifying the person's gender identity. However, especially in literary translations, such marking may shift the focus of the translation to gender in a way the source text did not require \citep[27]{nissen2002aspects}. 

In languages with a grammatical gender system, the same object or abstract concept can be assigned different genders and this can shape how it is perceived by speakers (see for example \cite{boroditsky2003sex}).\footnote{It should be noted that these results have been difficult to replicate \citep{mickan2014key}.} When translating, connotations conveyed by gender should be accordingly taken into account, especially in the case of metaphors and personifications. For instance, \textcite[237]{jakobson1959linguistic} notes how death is masculine in German (\textit{der Tod}), and depicted in books as an old man, but feminine in Russian, and therefore differently conceived. This issue can be tackled in translation by (i) looking for a synonym that has the same gender as the term to be translated; (ii) using a word from a third language that is also familiar to the readers of the translation; or (iii) providing a footnote \citep[29]{nissen2002aspects}. These strategies are also relevant for multimodal translation, such as the localisation of advertisements. Animals are often used to convey particular messages through the characteristics they are associated with because of their grammatical gender, e.g. using a tiger to evoke athleticism and therefore to convey the idea of speed in a petrol advertisement in German \citep[368--369]{di2020grammatical}.

Similarly, animals are often personified in fables or tales \citep[368]{di2020grammatical}. When this occurs, as in the Italian fable \textit{Il Ragno e l'Uva} (The Spider and the Grapes), the translation into notional gender languages is problematic. Even though in grammatical gender languages animals are assigned a gender, they are not perceived as having a male or female gender identity. For instance, even though the spider is masculine in Italian, the grammatical gender does not correspond to the animal's biological classification as male or female. However, translators in English are likely to mark gender since the use of the pronoun \textit{it} would weaken the anthropomorphic element of the text. This, in turn, adds a meaning for the target readers of the translation, i.e. the animal having a particular gender, that is not present in the source text \citep{di2020grammatical}. 

Another issue concerns social gender. Social gender is the property for which a word, usually a profession, is generally assigned the female or the male gender. This assignment is arbitrary and mainly based on stereotypical associations that vary across time and societies \citep[31]{nissen2002aspects}. For example, the word \textit{nurse} is often associated with femaleness. Social gender represents a problem when translating from notional gender languages into grammatical gender languages, since the nouns' and, therefore, a person's gender must be specified. This often leads to the source language input being translated differently across languages \citep[33--34]{nissen2002aspects}. For instance, a simple sentence such as ``My \textit{nurse} told me to take five pills a day'' requires that the gender of the nurse is specified when translating into Italian or German. In such cases, when there is no cue about gender, the assignment occurs based on social assumptions \citep[30]{nissen2002aspects}.

More generally, translation from notional gender languages into grammatical gender languages can force translators to opt for a gender \citep[367]{di2020grammatical}. In literary texts, authors can avoid gender markers so as to create a mysterious atmosphere, for example in Edgar Allan Poe's \textit{The Tell-Tale Heart} (\cite*{poe1843tell}), where it is impossible to define the narrator's gender as they speak in the first person. \textcite[369]{di2020grammatical} note how the Italian translator opted for masculine forms of nouns, adjectives, and participles, thus choosing a particular gender for the narrator. 

Finally, \textcite[369]{di2020grammatical} claim that the translation of occupations and/or institutional roles is related not only to semantics but also to ideological questions. For example, feminine forms of high-level professions and institutional roles in Italian are not particularly widespread although they are semantically and grammatically correct (see \cite{gheno2019femminili}). In Italy, their use is often associated with ideological matters rather than simple linguistic correctness \citep[112]{gheno2019femminili}. Similar claims can also be made for GFL. Its use, for example in a translation in order to replace masculine generics, may be seen as a political and/or ideological statement rather than a simple matter of equality. In such a situation, translators can help build bridges between marginalised communities and their clients but they could also be denied an assignment because of their willingness to opt for GFL strategies.

The aforementioned issues show how gender in translation can be a complex and ambiguous phenomenon and every solution adopted has different connotations and implications. There are no clear, one-size-fits-all solutions on how to handle gender appropriately and translators should, while being creative and experimenting with language, consider a wealth of co-textual and contextual factors. The translation of gender represents a challenge for MT as well. This is thoroughly discussed in Subsection \ref{sub:biasMT}.  

\section{Sex/Gender: Introduction} \label{sec:sex/gender}

This section introduces some key concepts that permeate the field of gender studies. It focuses on sex and gender, more specifically their meanings and interconnections, their importance, and the role of society in shaping assumptions concerning them. Attention is paid to how the concepts of sex and gender are, at least partially, socially constructed and still conceived as binary, i.e. male/female, man/woman. Since GFL in the present work is intended as a language that includes all or no genders, an overview of non-binary\footnote{The term non-binary is used here to refer to all sexes and genders beyond male and female, also including intersex people.} gender identities is provided. 

%----------------------------------------------------------------------------------------
%	SUBSECTION 1
%----------------------------------------------------------------------------------------

\subsection{Sex vs gender}

Gender identity must be indicated to access numerous goods and services, e.g. to open a bank account or to enrol at a university \citep{fae2016non}. What is asked in forms or on websites, however, is not always clear since the terms sex and gender are often used synonymously in different contexts (\cite[194f]{van2017biological}, \cite[105]{acanfora2021in}). 
%Sexuality as well revolves around gender. \textcite[1177]{van2015beyond} claims that gender provides the foundation to label sexuality with terms such as heterosexual, homosexual, bisexual, etc. This is indeed defined by a person's gender on the one hand and the genders of the people they find sexually attractive on the other. For instance, a man sexually interested in women is defined as heterosexual, while if he is interested in men, he is defined as homosexual. The researcher proposes to address the complexity of sexuality through \textit{Sexual Configurations Theory (SCT)}, affirming that many other elements define it (see Fig. \ref{fig:SCT Theory} for a simplified overview), gender being only one of them. She also claims that it is difficult to distinguish gender from sex in this regard, wondering if a person is attracted by physical features, social identities, interactions or other elements (\cite[1177]{van2015beyond}). Such considerations can be used as a starting point to introduce the discourse about sex and gender. 
%-----------------IMAGE begins here------------------
%\begin{figure}[ht]
%\centering
%\includegraphics[width=1.0\textwidth]{figures/SCT.pdf}
%\caption[SCT Theory]{Simplified overview of the defining elements of sexuality according to SCT Theory %(\cite[12--13]{iantaffi2018mapping})}
%\label{fig:SCT Theory}
%\end{figure}
%-----------------IMAGE ends here-----------------
%First, the terms sex and gender are used often interchangeably (\cite[105]{acanfora2021in}, \cite[194f]{van2017biological}). 
It is generally assumed that a person with certain physical characteristics (i.e. sex) will identify as a certain gender. For instance, it is taken for granted that a person born with female genitalia will identify as woman (\cite[95]{richards2016non}, \cite[1]{morgenroth2020effects}). In such a view, gender follows from sex. It is very difficult to find an uncontentious and clear definition of both terms as well as to understand the difference between them (\cite{lester2018trans}). 

Sex is generally described as referring to biological characteristics, such as chromosomal makeup, hormones, and primary and secondary sexual characteristics, whereas gender is defined as a historical, cultural, and social construct that refers to roles and behaviours associated with the equally socially constructed concepts of masculinity and femininity (\cite[122]{barker2018rewriting}, \cite*[32]{barker2018psychology}, \cite[56]{barker2019life}). In the heterosexist Western society, both sex and gender are thought to be binary and this binarism not only dictates how sex and gender are linked -- with the latter supposedly following from the former -- but also how they influence sexuality (\cite[23--24]{butler1990gender}, \cite[123]{barker2018rewriting}). As shown in Figure \ref{fig:sex, gender, sexuality}, it is generally assumed that a person is born having one sex (either male or female) and that they will congruently grow up being a masculine man if male or a feminine woman if female. According to their sex/gender, they will also be attracted to a person of the opposite sex/gender \citep[123]{barker2018rewriting}. This is, however, an essentialist view of gender.

%---------IMAGE begins here----------
\begin{figure}[ht]
\centering
\includegraphics[width=1.0\textwidth]{figures/sex, gender, sexuality.pdf}
\caption[Sex, gender, sexuality]{Interconnections among sex, gender and sexuality in a heteronormative society \citep[123]{barker2018rewriting}}
\label{fig:sex, gender, sexuality}
\end{figure}
%------------IMAGE ends here------------

%Media, gender and identity: An introduction andrebbe citato ma non trovo il libro in formato dgtl prcmdnn

%-----------------------------------
%	SUBSECTION 2
%-----------------------------------
\subsection{Defining sex} \label{Sex and Biology: an overview of the complexity}

The term sex, often interpreted as biological sex, is used to refer to biological features. However, when referring to a person's sex, the term sex assigned at birth is preferred and even recommended (\cite{fausto2012sex}, \cite*{fausto2018why}, \cite[24]{prona2016mein}, \cite[95]{richards2016non}, \cite{serano2017transgender}, \cite[30]{barker2018psychology}, \cite{lester2018trans}, \cite[57]{barker2019life}, and see also recommendations by \cite{apagenderenglish}). At birth, sex is assigned based on the sole inspection of the newborn's genitalia provided that they are dimorphic, i.e. they can be clearly categorised as female (clitoris) or male (penis). This has historically erased the identity of intersex people.\footnote{Intersex people, however, can have perfectly dimorphic genitalia. This is further explained in the present section.} Individuals who were born with mixed genitalia have usually undergone a surgical operation to fit dimorphism \citep{fausto2012sex,fausto2018why,fausto2020sexing,prona2016mein,richards2016non,barker2018rewriting,lester2018trans,barker2019life}. Three different arguments can be advanced for preferring the expression sex assigned at birth over the simplistic term sex when referring to a person:
\begin{enumerate}
\setlength\itemsep{-0.5em}
    \item If human biology in its real complexity is taken into account, then five layers of sex should be examined. They are not always binary and sometimes also develop independently of one another \citep{fausto1993five,fausto2000five,fausto2012sex,fausto2018why}.
    \item Science is always produced and interpreted within a specific sociocultural context. In this case, biologists produced or interpreted facts, namely the distinction between male and female, based on their assumptions on the gender binary (see \cite{sterling2020science}). Sex can therefore be considered in part as socially constructed as gender (\cite[10--11]{butler1990gender}, \cite[199]{fausto2016critiques}, \cite{serano2017transgender}).
    \item Biological features are actually malleable so that if a person chooses to undergo hormonal and/or medical treatment, their biology as well as their primary and secondary sex characteristics change \citep{serano2017transgender,lester2018trans}.
\end{enumerate}


Fausto-Sterling (\cite*[66]{fausto2012sex}, \cite*[190]{fausto2016critiques}) argues that the human body is extremely complex and many traits such as voice pitch or breast size are not dimorphic. Rather, they overlap and make it difficult to provide precise answers about sex differences. As regards sex and its development, the biologist distinguishes five layers (\cite[3]{fausto2012sex}, \cite*{fausto2018why}). The first layer of sex develops at conception and (i) is the chromosomal sex: a sperm bearing X or Y chromosomes normally fuses with an egg, which bears X chromosomes. From this, the common misconception originates that a person with XX chromosomes is female and a person with XY is male. Other combinations, such as XXY, XYY or even XO, are possible (\cite[25]{fausto2012sex}, \cite*{fausto2018why}). After eight to twelve weeks the embryo develops a pre-gonadal structure that can develop either in a male or female direction, but unexpected developments can also occur, which is the (ii) foetal gonadal sex. At this stage, primary sex characteristics (testes and ovaries) are determined. Hormones developed by the embryo, this is the (iii) foetal hormonal sex, push the development of testes and ovaries further, determining their differentiation and leading to the formation of (iv) the internal reproductive sex, i.e. uterus, cervix and so forth in females or prostate, epididymis and so forth in males. Finally, in the fourth month, (v) the external genital sex is developed (\cite[5]{fausto2012sex}, \cite*{fausto2018why}). Fausto-Sterling (\cite*[20]{fausto1993five}, \cite*[20]{fausto2000five}, \cite*[11]{fausto2012sex}, \cite*[190]{fausto2016critiques}, \cite*{fausto2018why}) has more than once claimed that different combinations at every layer of sex are possible and it would be better to refer to sex as a continuum rather than as binary. This complexity is usually reduced to an essentialist model and \textcite[190]{fausto2016critiques} underlines that the term sex is used in various different ways in both science and everyday language. To summarise, if sex is used to account for psychological, physiological or sexual behaviours and differences, it is not a useful and precise concept.

Discoveries in the field of sex development and how they were interpreted led to different theories. \textcite[10--11]{butler1990gender}, \textcite{fine2017testosterone} and Fausto-Sterling (\cite*{sterling2020science}, \cite*{fausto2020sexing}) claim that the production as well as the interpretation of scientific facts is always influenced by cultural assumptions. For instance, gender is believed to be binary and to follow from sex as well as alleged (stereotypical) differences between men and women. The scholars conclude that sex too is, at least in part, socially constructed like gender. \textcite{fine2017testosterone} analysed extensive research in the fields of biology, psychology, and sociology. She noted, for example, how testosterone and other hormone levels were once thought to be genetically determined as well as fixed, leading to the use of such arguments to support claims on the nature of the different roles of men and women in society. The human brain too was believed to differ based on biological sex, but dimorphism is actually rare (see \cite{fine2017testosterone} and \cite[172]{hyde2019future} for an overview of the existing literature on the topic). Nowadays, scholars are unanimous on the position that biological features are not fixed, but are rather influenced by social and environmental conditions. For example, fathers, and especially those who invest more time in caregiving, show decreased levels of testosterone (\cite{gettler2011longitudinal}). Such levels increase in female actors who enact power roles, thus reproducing stereotyped roles (\cite{van2015effects}). Differences in hormone levels or human brains are a product of sex/gender differences rather than a cause thereof since we live in a highly gendered society (\cite[200]{fausto2016critiques}, \cite*{sterling2020science}, \cite{fine2017testosterone}, \cite[142]{barker2018rewriting}, \cite*[36]{barker2018psychology}, \cite{lester2018trans}, \cite[174]{hyde2019future}, \cite[58]{barker2019life}). This shows how research is sometimes biased towards social assumptions that affect how experiments are carried out and results interpreted (see \cite{fine2017testosterone} and \cite{hyde2019future} for an overview in the field of sex/gender-related differences as well as \cite{fausto2020sexing} for a very detailed account of how changing societal assumptions shaped interpretations of discoveries in the field of sex development).

In addition, transgender activists such as \textcite{serano2017transgender} and \textcite{lester2018trans} explain how biological makeup, primary (e.g. genitalia) and secondary characteristics (e.g. facial hair) can change through hormonal or surgical intervention. This holds also for cisgender people (see Subsection \ref{identities} for a definition of this term), who may have their breasts removed after cancer or take the birth control pill, for example. Lester notes:
\begin{quote}
    The body of a trans woman who has pursued hormone therapy and surgery might combine XY chromosomes, higher levels of estrogen and progesterone with concurrent lower levels of/lower sensitivity to androgens, no testes, a prostate, a vulva and vagina, little body hair, no facial hair, breasts and curves. When looking at all the different parts of her physical makeup, what counts as “biological sex”? All of these categories are sexed, and all are “biological.” \citep[60]{lester2018trans}
\end{quote}

The process of sex assignment can be further detailed by looking at intersex movements. A newborn with a penis is not only assigned male at birth but is also recognised as such by society and law (\cite{serano2017transgender}, \cite[56]{barker2019life}). However, if in the past one's reproductive anatomy did not fit dimorphism, i.e. the external genitalia could not be clearly labelled as clitoris or penis, genital surgery was readily performed for genitals to fit a perfect dimorphism. The children were then raised according to the surgically assigned gender (\cite[19--20]{fausto2000five}, \cite*{fausto2020sexing}), thus systematically erasing the identity of intersex people to make them conform to binary social and cultural expectations (\cite[57--58]{barker2019life}). As intersex movements rose, intersex people started to advocate for more ethical treatments that should include surgery only as a means to save lives; they also demanded that intersexuality be defined as ``normal'' (\cite[21]{fausto2000five}, \cite[57--58]{barker2019life}). 

Some intersex people may not find out about their intersex status until adolescence; other may never find out at all. For example, some XX babies can be born with a penis, but as long as the external genitalia are clearly distinguishable, no further test is carried out at birth (\cite{fausto2018why}, \cite[57]{barker2019life}). Also, people have their chromosomal makeup examined quite rarely (\cite{lester2018trans}). This means that they could be assigned the wrong sex at birth. Acknowledging all the aforementioned points is a step towards dismantling transphobic arguments. Fausto-Sterling claims, describing the treatments given to intersex people, that:
\begin{quote}
    labeling someone a man or a woman is a social decision. We may use scientific knowledge to help us make the decision, but only our beliefs about gender -- not science -- can define our sex. \citep[3]{sterling2020science}
\end{quote}

This quotation confirms that labels, as explained in Subsection \ref{sec:languageculture}, are influenced by culture and society. Therefore, sex is categorised differently in other, non-Western cultures. The question of how sex and/or gender identities are represented is a crucial matter not only for gender studies but also for TS \citep{spurlin2014gender}, highlighting its interdisciplinary character.

%Effects of gendered behavior on testosterone in women and men (van Anders et al. 2015)
%Gettler, L., McDade, T., Feranil, A., & Kuzawa, C. (2011). Longitudinal evidence that fatherhood decreases testosterone in human males. Proceedings of the National Academy of Sciences, 108(39), 16194–16199.

Finally, \textcite{serano2017transgender} argues that efforts to discriminate against trans people, especially trans women, resort to the alleged distinction between sex, hence biology, and gender, hence social construction. According to such distinction, transphobic people would argue that a trans woman is biologically male. Similar arguments are used by Trans-Exclusionary Radical Feminists (TERF) to exclude trans women from feminism (see \cite{lester2018trans}). Understanding how the category of sex is interpreted by social assumptions, how people are required and supposed to fit a perfect dimorphism and how sex is diverse helps dismantle such arguments. To this end, transgender movements also ask that their opinion be considered by the scientific community and that trans people be included in the process of diagnosis or classifications about transsexuality or gender dysphoria.\footnote{The negative feeling that a person's body and experiences do not fit with how they perceive themselves \citep[65]{barker2019life}.} They claim that medical language pathologises them and that science is often imprecise, delivering ill-founded conclusions \citep{sterling2020science}. In this regard, \textcite[76--77]{acanfora2021in} affirms that diagnoses are made in a pathologising, static, normalising, and fixing language that is then often also used in everyday life, thus the diagnosed person tends to identify with their diagnosis. 

%--------------------------------------
%	SUBSECTION 3
%-------------------------------------

\subsection{Defining gender} \label{Gender}

Before 1955, the term \textit{gender} was almost exclusively used to refer to grammatical gender (see Subsection \ref{sub:language/gender}) \citep[89]{haig2004inexorable}. In 1955, \textcite{money1955hermaphroditism} introduced the concept of gender role to refer to social behaviours and characteristics that are associated with masculinity and femininity \citep[89]{haig2004inexorable}. From the second half of the 1960s, feminists then began to use the term \textit{gender} to account for differences in men's and women's societal roles and opportunities that were the results of social processes and could not be ascribed to biology \citep[5]{von2016translation}. Nowadays, in public discourse, gender is still often used as a synonym for sex, and there are at least two views of gender: essentialist and constructionist \citep[106]{acanfora2021in}.

The essentialist view of gender posits that some gender features are innate, universal, and biologically determined \citep[106]{acanfora2021in}. The constructionist view posits that gender is a historical, cultural, and social construct,
tied to the notions of masculinity and femininity (\cite[122]{barker2018rewriting}, \cite[32]{barker2018psychology}, \cite[56]{barker2019life}, \cite[106--107]{acanfora2021in}). A third view was proposed by \textcite[25]{butler1990gender} who introduced the concept of gender performativity. The concept is based on the speech act theory proposed by \textcite{austin1975things}, who describes explicit primary performative utterances, i.e. utterances that simultaneously perform an act and describe its performance, e.g. ``I promise I will call you tomorrow''. The actions in the utterances are therefore produced through their description. Similarly, through discourse, gender performatively constitutes itself and it is treated as something that already exists. Gender is therefore performed by being enacted and, at the same time, its categories are (re)created and presented as natural \citep[25]{butler1990gender}. In this process of (re)creation of gender, language and translation play an essential role because they can uphold or challenge societal views (\cite{simon1996gender,sprachhandeln2015tun,fairclough2015language}) as discussed in Section \ref{TSrelevance} and Subsection \ref{languapepowerdiscrimination}.


In Western societies, gender was (and often still is) believed to be binary and to follow from sex (\cite[95]{richards2016non}, \cite[116--117]{matsuno2017non}, \cite[122]{barker2018rewriting}, \cite[28]{barker2018psychology}, \cite{lester2018trans}, \cite[56--57]{barker2019life}). \textcite[174]{butler1990gender} tried to dismantle such a view, describing different forms of what she defined as \textit{gender troubles}, i.e. people that challenge the performative binary gender system, such as drag queens. The binary system is still believed to be a valid model, with the consequence that men and women are encouraged to take on specific roles and act in specific ways.

For instance, \textcite[114]{heilman2012gender} defines gender stereotypes as ``generalizations about attributes of men and women'' that include their main features as well as inclinations for particular roles and/or responsibilities. They can be divided into descriptive and prescriptive stereotypes. The former describe what men and women are like; the latter, what they should be like. Stereotypes thus dictate how people have to act and devalue whoever does not fit them \citep[114]{heilman2012gender}. As shown in Table \ref{tab:genderstereotypes}, \textcite{heilman2012gender} argues that women are described as having communal traits, and men agentic ones. 


%-------------------------------------------TABLE BEGINS HERE---------------------------------------------
\begin{table}
    \caption{Gender stereotypes according to \textcite[115]{heilman2012gender}}
    \begin{tabularx}{\textwidth}{QQ}
    \lsptoprule
         \textbf{Agency} & \textbf{Communality}  \\
         \midrule
         achievement-orientation (competent, ambitious, task-focused) & concern for others (kind, caring, considerate)\\
         inclination to take charge (assertive, dominant, forceful) & affiliative tendencies (warm, friendly, collaborative)\\
         autonomy (independent, self-reliant, decisive) & deference (obedient, respectful, self-effacing)\\
         rationality (analytical, logical, objective) & emotional sensitivity (perceptive, intuitive, understanding)\\
         \lspbottomrule
    \end{tabularx}
    \label{tab:genderstereotypes}
\end{table}
%---------------------------------------------TABLE ENDS HERE----------------------------------------------

Gender stereotypes are often used by the media, especially for marketing reasons, to get the attention of the public \citep[124]{barker2018rewriting}. They are also perceived as rules and negatively affect men, women, and non-binary people \citep[132]{barker2018rewriting}. Men have to be rational, dominant, and protective of other people such as their partners or children. They must not be intimate with other men and feminine men are discriminated against since femininity is despised. Women are often defined by their relationships with other people, especially with men, discriminated against in the workplace, or seen as mere caregivers or sex objects (\cite{heilman2012gender}, \cite{grover2015second}, \cite[132]{barker2018rewriting}, \cite{murgia2021stai}).
Non-binary people are subject to even more discrimination: they are misgendered, denied access to public facilities such as lavatories, verbally and physically harassed or their existence is even denied (\cite{herman2013gendered}, \cite{mclemore2015experiences}, \cite*{mclemore2018minority}, \cite[118]{matsuno2017non}, \cite{lester2018trans}). Gender roles and stereotypes originating from a binary and outdated understanding of gender thus prove harmful to many people. The discourse around gender stereotypes is relevant for TS as well, e.g. because marketing advertisements playing with them are translated into other languages \citep{di2020grammatical}, once again showing the social relevance and impact of language and TS. 

%---------IMAGE begins here----------
\begin{figure}[ht]
\centering
\includegraphics[width=0.7\linewidth]{figures/biopsychosocial.pdf}
\caption[Biopsychosocial processes]{Biopsychosocial processes and their interaction
(\cite[36]{barker2018psychology})}
\label{fig:biopsychosocial}
\end{figure}
%------------IMAGE ends here------------

A newer understanding of gender consists in defining it as a biopsychosocial construct, hence biological, psychological, and social aspects interacting with one another, as shown in Figure \ref{fig:biopsychosocial} (\cite[99f]{barker2018rewriting}, \cite*[35]{barker2018psychology}, \cite[58]{barker2019life}). Biological aspects may involve brain functionality and hormone production as well as their interaction. Psychological aspects may involve the way people experience gender, e.g. how they identify, and how they think about gender. Social aspects involve how gender is enacted in a particular historical, geographical, and cultural context, hence expectations about genders, gender roles and norms, e.g. how one has to dress, which activities one takes up, etc. \citep[58]{barker2019life}. These three aspects are all interrelated and cannot be considered separately as they influence one another. Thus, gender is a construct, but it influences biology and also other tangible aspects of human life. Since the binary, static, and fixed gender system is deeply rooted in Western cultures, people have to locate themselves in relation to it. They are deeply influenced by it, e.g. in the way they dress, act and speak, and cannot completely abandon this model (\cite[81]{barker2019life}, \cite[108f]{acanfora2021in}).
Besides, when talking about a person’s gender, reference could be made to their: (i) gender identity: the gender they identify as, like woman, man, genderfluid, non-binary, genderless, etc.; (ii) gender role: how feminine and/or masculine they act according to gender norms; (iii) gender expression: how they express their gender through clothes, hairstyle, etc.; (iv) gender experience: how they face the world based on their gender (\cite[129]{barker2018rewriting}, \cite[60]{barker2019life}). Finally, gender also interacts with sex. 



%--------------------------------------
%	SUBSECTION 4
%-------------------------------------

\subsection{Gender identities: Cisgender, Transgender, and other terms} \label{identities}

In Subsection \ref{Sex and Biology: an overview of the complexity}, it was argued that sex is assigned at birth based on certain physical characteristics, namely the form of the external genitalia:
\begin{quote}
If people are assigned a certain sex at birth and identify with a gender that is considered congruent with that sex assigned at birth culturally, they are cisgender—that is their sex assigned at birth and gender identity are on the same side. \citep[59]{barker2019life}
\end{quote}
When people identify with a gender that is not congruent with the sex assigned at birth, they might be considered transgender \citep[60]{barker2019life}.
More specifically, by the fourth month of pregnancy, genital sex is developed and therefore sex is often assigned before birth through inspection of the baby's genitalia. A process of genderisation begins: medical professionals and parents expect the baby to identify as the gender which is socially held as congruent with their genital sex, and parents buy clothes and toys accordingly, e.g. dolls, skirts, or other pink clothes for female babies, cars and blue clothes for male children. Through clothes and colours, children's sex assigned at birth is made visible so that adults can make assumptions about their genitalia (\cite[7]{fausto2012sex}, \cite*[195]{fausto2016critiques}, \cite[107]{acanfora2021in}). Hence, when a person is assigned a sex before or at birth, e.g. female, they are also assigned a specific gender based on social expectations, e.g. woman. If the newborn identifies as the gender assigned at birth, they might be defined as cisgender. Otherwise, they might be defined as transgender. 

This distinction may seem straightforward; however, it is also a binary concept, which is problematic. For example, according to \textcite[72]{barker2019life}, cross-dressers, such as drag queens, who perform by temporarily impersonating another gender, would fall under neither the cisgender nor the transgender umbrella. On the contrary, \textcite{lester2018trans} argues that cross-dressers fall under the transgender umbrella. Even though there is no unanimity, the terminology presented in this subsection helps move away from the category of biological sex. 

Moreover, there are intersex people, i.e. people who are born with variations in sex characteristics. As explained in Subsection \ref{Sex and Biology: an overview of the complexity}, these include chromosomes, gonads, hormones, and genitals. In this case, the use of the term cisgender or transgender may be controversial. Since the use of labels, such as transgender, is also a question of personal preference, one should always ask for the terms people identify with and use them.

That said, it is useful to define the term non-binary. Non-binary is both a term describing a gender identity and an umbrella term that includes all genders other than man and woman (\cite[96]{richards2016non}, \cite[150]{barker2018rewriting}, \cite*[28]{barker2018psychology}, \cite{lester2018trans}, \cite[61--62]{barker2019life}). It can be abbreviated to \textit{enby} or NB, but some people also use the term genderqueer (\cite[96]{richards2016non}, \cite[61--62]{barker2019life}). Labels for gender identities are numerous and include different experiences of gender (\cite[96]{richards2016non}, \cite[150]{barker2018rewriting}, \cite[61]{barker2019life}): 
\begin{itemize}
\setlength\itemsep{-0.5em}
    \item No sense of gender (e.g. agender, genderless, etc.).
    \item Muted/partial sense of gender (e.g. demi boy or demi girl).
    \item Both binary genders, i.e. male and female (e.g. androgynous, mixed gender, etc.).
    \item Changing over time (e.g. genderfluid, bigender).
    \item A third gender (e.g. third gender, other gender).
    \item Multiple genders (e.g. pangender).
    \item People who challenge the ontology of gender (e.g. genderqueer, genderfuck, etc.).\footnote{Note that this list does not claim to be exhaustive and representative of all gender identities and experiences. Also, note that some terms may overlap and one must be aware of people's identified labels.}
\end{itemize}
At this point, it is useful to sum up the different meanings of the terms cisgender and transgender. 
Cisgender people's gender identity is aligned with the gender that is socially expected based on their sex assigned at birth. They can be either (cis-)men or (cis-)women. Transgender people's gender identity is not aligned with the gender that is socially expected based on their sex assigned at birth. They can have a binary gender, hence being (trans-)men or (trans-)women. The former were assigned female at birth and identify as men, while the latter were assigned male at birth and identify as women. Transgender people may also be non-binary, which is an umbrella term for gender identities outside the binary. However, not every non-binary person defines themselves as transgender. Finally, people who cross-dress may or may not be considered transgender (see Figure \ref{fig:gender umbrella} for a simplified illustration of such umbrella terms) (\cite[96]{richards2016non}, \cite[150]{barker2018rewriting}, \cite{lester2018trans}, \cite[61]{barker2019life}, \cite[109]{acanfora2021in}). Transgender people may or may not transition -- this is a long and complex process to move from living in the sex assigned at birth to living in one's gender identity. The transition may be social, e.g. it may involve coming out as transgender, changing one's name, pronouns and documents, and/or physical, e.g. it may consist of hormone treatments and different types of surgery to align the body with one's perception of how this should be. Transition and the representation of transgender people are also a research object of queer translation: for instance, \textcite{rose2016keeping} examines how to preserve the gender identity of trans people in the translation of two memoirs from French, where both masculine and feminine gender markings are used in reference to the same transgender person, into English.

%-----------------IMAGE begins here------------------
\begin{figure}[!ht]
\centering
\includegraphics[width=0.9\textwidth]{figures/genderumbrella.pdf}
\caption[Gender umbrellas]{Simplified illustration of the gender umbrellas with gender identity terms}
\label{fig:gender umbrella}
\end{figure}
%-----------------IMAGE ends here-----------------

%-----------------IMAGE begins here------------------
\begin{figure}[!ht]
\centering
\includegraphics[width=0.7\textwidth]{figures/genderunicorn.pdf}
\caption[Gender Unicorn]{Simplified illustration of sex/gender identity, sexual and romantic orientation \citep{trans2015unicorn}}
\label{fig:gender unicorn}
\end{figure}
%-----------------IMAGE ends here-----------------

In Subsection \ref{Gender}, it was argued that when talking about gender, reference could be made to different aspects, such as gender identity, role, expression, and experience (\cite[60]{barker2019life}). In this subsection, labels for different gender identities were explained. The interaction of different aspects of gender as well as of sex and sexuality has been summarised and illustrated in some models such as in Figure \ref{fig:gender unicorn} by \textcite{trans2015unicorn}. This is an adjustment of an outdated version of the model proposed by \textcite{killermann2017genderbread}\footnote{Even though \textcite{killermann2017genderbread} updated his model in 2017, the re-elaboration proposed by \textcite{trans2015unicorn} is to be preferred as it is more terminologically precise.} to more accurately depict the differences between gender (identity and expression), sex assigned at birth, and sexuality. Note that physical and emotional attraction are regarded as separate, so a person could be emotionally attracted to more genders but physically attracted to one gender only (see for example \cite[1200f]{van2015beyond}, \cite[110]{barker2018rewriting}, \cite[40]{barker2019life}). Such models, however, are extremely simplified and not representative of our gender and sexual identities.


%--------------------------------------
%	SUBSECTION 5
%-------------------------------------

\subsection{Beyond sex and gender}
\textcite{van2015beyond} and \textcite{iantaffi2018mapping} proposed Sexual Configurations Theory (SCT) because, they maintain, the classical model of sexual orientation poses several problems: (i) it is based on a binary and confused conception of gender (and sex) and it is not clear what is meant by the term; (ii) other aspects of sexuality are ignored, gender being only one of them, e.g. levels of sexual attraction, number of partners involved in the sexual act, and the roles played (see Figure \ref{fig:SCT Theory} for a simplified overview); (iii) it is a static and fixed model that does not account for sexual fluidity and the existing variety of experiences of sexuality \citep[1178f]{van2015beyond}. 

%-----------------IMAGE begins here------------------
\begin{figure}[!ht]
\centering
\includegraphics[width=0.8\textwidth]{figures/SCT.pdf}
\caption[SCT]{Simplified overview of the defining elements of sexuality according to SCT \citep[12--13]{iantaffi2018mapping}}
\label{fig:SCT Theory}
\end{figure}
%-----------------IMAGE ends here----------------

\textcite[1178]{van2015beyond} stated that she conceived ``Sexual Configurations Theory (SCT) as a way to address the complexities of actual people’s sexualities''. In the context of this book, it is important to underline that van Anders chose to use the term gender/sex, as she claims that sex and gender cannot always be easily distinguished and some phenomena cannot be sorted into one or the other, e.g. the question is whether hormones, which are a biological feature whose level nevertheless changes based on the social context, are related to sex or gender (\cite[1181--1182]{van2015beyond}, \cite{lester2018trans}, \cite[59]{barker2019life}). 

%In addition, \textcite[1189]{van2015beyond} offers an interesting analysis of the multi-faceted aspects of gender/sex in relation to sexual attraction that can be very useful to understand the complexity and the variety of gender/sex identities (see Figure \ref{fig:gender/sex}).

%\begin{figure}[ht]
%\centering
%\includegraphics[width=0.9\textwidth]{figures/gender_sex.pdf}
%\caption[Gender/Sex]{The complexity of gender/sex identities (\cite[14]{iantaffi2018mapping})}
%\label{fig:gender/sex}
%\end{figure}

%The dimension of gender/sex type is represented by a continuous ring that is not closed in the middle: it describes attraction to binary gender/sexes. Attraction is always conceptualised as continuous. If a person is attracted to binary gender/sexes, they could be attracted to both men and women but to different extents. The middle of the ring represents attraction to individuals outside the binary: such non-binary gender/sex level refers to people that challenge the binary, it can span from genderqueer people to intersex individuals and many others. At the very center there are people who completely challenge gender norms but one could challenge more men- or women-specific norms if it was closer to the straight repeated lines, for example. Attraction can also span from a specific gender/sex, to both binary gender/sex, more genders or to all genders. Even though the scheme presented in Fig. \ref{fig:gender/sex} was thought to represent possible interaction between gender/sex and sexuality, I would argue that it also helps understand the variety and fluidity of gender/sex identities. For instance, a person could identify as man, bigender, genderfluid or pangender. Moreover, gender/sex strength refers to how gender/sex is important to sexuality from 100\% to 0\% and, also in this case, this could be compared to agender or genderless individuals, i.e. people who do not have any sense of gender (\cite[1189]{van2015beyond}). Finally, to acknowledge the complexity of sex and gender as well as their intercorrelations means to understand the costs of the gender/sex binary.

Once again, it should be noted that the gender, and also sex, binary is the cause of different types of discrimination. The existence of individuals whose identities or bodies fall outside the binary, as well as that of people whose gender identity is binary but not coincident with the sex assigned at birth, is hardly acknowledged. They are discriminated against by being misgendered, denied access to public facilities such as lavatories, and verbally and physically harassed (\cite{herman2013gendered}, \cite{mclemore2015experiences}, \cite*{mclemore2018minority}, \cite[118]{matsuno2017non}, \cite{lester2018trans}, \cite[184]{hyde2019future}). 
Cis-men and cis-women are then believed to be very different from one another and have to fit old gender roles and stereotypes that affect them in various ways,\footnote{This is also true for transgender people: for them to access medical treatments for transitioning, they are required to fit the stereotypical characteristics of femininity or masculinity. This is especially discriminating for transgender non-binary people as they could be more easily denied hormone therapy and/or surgery (\cite[28]{prona2016mein}, \cite{lester2018trans}).} for example preventing women's access to leadership positions (\cite{grover2015second}, \cite[132]{barker2018rewriting}, \cite[183--184]{hyde2019future}). Such beliefs negatively affect one's self-perception, so that one thinks about people's identity in a stereotyped manner (\cite[184]{hyde2019future}). This also shapes legal and social policies. Single-sex schooling was introduced based on alleged differences between male and female brains. Hormone levels, and in the past also genetic testing and the examination of genitalia, were and are used to regulate women's access to sports competitions \citep{sterling2020science}. Sex or gender -- it is not always clear which criteria the policies employ -- are used as a basis to manage access to public spaces (\cite{van2017biological}, \cite[184]{hyde2019future}, see \cite{sterling2020science} for some accounts of the history of sports competitions and the discrimination of intersex people). Finally, the gender binary has proven to be an obstacle to advances in research. Since gender/sex includes a wider variety of phenomena, the use of simple variables, such as female/male or woman/man, fails to acknowledge the causes (e.g. hormones, gender roles expectations, socioeconomic status) of differences for example in depression rates among women, men and individuals of other genders \citep[184]{hyde2019future}. Gender and sex are hence intertwined, biological and social/cultural processes interacting and influencing one another and causing a multitude of differences. 


%--------------------------------------NEW PAGE-------------------------------------------
%--------------------------------------
%	SECTION 2
%-------------------------------------
\section{The societal impact of gender in language} \label{sec:powerlanguage}

The concept of gender was introduced in the previous section with reference to its biopsychosocial nature. Gender also interacts with language. It is expressed through linguistic structures which reflect and create assumptions about it \citep{stahlberg2007representation}. The following sections deliver insights into the interplay of language, society, and gender, providing a better understanding as to why gender issues have received great attention in the field of TS as well. Languages are the result of cultural conventions, thus they express the same events or concepts in different ways. Such differences, in the case of gender, not only complicate the translation process as shown in this book (see Chapter \ref{results}) but they also may impact people's perceptions. Since language -- and therefore translation \citep{wolf2007sociology} -- is socially situated, the reverse also happens: societal norms and power relations, such as sexism, are (re)produced through language and translation \citep{simon1996gender,von2016translation}. This means that language can also be used to foster change. The chapter concludes with a discussion of empirical research that shows how the use of GFL positively impacts gender equality. 

%--------------------------------------
%	SUBSECTION 1
%-------------------------------------
\subsection{Language and thought} \label{sec:languageculture}
Language is a cultural convention \citep[9f]{deutscher2010through} and it can be used to express abstract and concrete concepts. Concepts are named differently across languages, e.g. happiness is \textit{felicità} in Italian but \textit{felicidad} in Spanish. The distinction among concepts may be based on natural differences. For example, the difference among animals, such as cats and dogs, is the same in every language. In other cases, the distinction among concepts may be the result of cultural conventions. For instance, the differentiation among body parts, such as arm and hand, is not the same in every language and culture \citep[13]{deutscher2010through}. 

Languages differ considerably in their structures as well. In Tagalog, a language spoken in the Philippines, different pronouns are used for \textit{we}, depending on the different meanings, i.e. (i) you and I; (ii) someone else and I; and (iii) someone else and I but not you \citep[15]{deutscher2010through}. Other languages, such as Mian, spoken in Papua New Guinea, use different tenses specifying whether a past event took place in an immediate, recent or distant past, whereas others, such as Indonesian, do not \citep[64]{boroditsky2001does}. These differences lead people to use different linguistic structures and hence can be assumed to think differently of the same event to express it \citep[64]{boroditsky2011language}. 

This idea originates from \textcite{whorf1956language}. He was a linguist who posited that language affects cognitive reasoning (see also \cite[1014]{pederson2007cognitive}). Such a notion became famous as the Sapir-Whorf hypothesis. It is also known as the theory of linguistic relativity and has long been debated in the field of cognitive linguistics (see \cite{pederson2007cognitive} for an overview). \textcite[72]{kramer2016feminist} argues that a version of this theory, namely a strong linguistic relativism or linguistic determinism, initially became popular. According to this idea, language influences logical reasoning and people are not able to understand concepts foreign to their language and perceive reality in a completely different manner (\cite[21--22]{deutscher2010through}, \cite[72]{kramer2016feminist}). 

This version is nowadays rejected, with linguists putting forward a moderate version (see for example \cite{levinson1994introduction,boroditsky2001does,boroditsky2003sex,phillips2003can,casasanto2008s,prewitt2012gendering,santacreu2013female}), investigating whether differences or similarities in linguistic categories such as time, space, and gender are correlated with cognitive and behavioural differences \citep[72]{kramer2016feminist}. For instance, time is conceived differently based on linguistic metaphors used to describe it, e.g. when saying ``move a meeting forward'' instead of \textit{up} or \textit{right} \citep[7]{boroditsky2001does}. Or again, speakers of numerous aboriginal languages like Guugu Yimithirr use only absolute directions to refer to the position of objects and people and develop a better sense of orientation than speakers of languages like English that use relative directions \citep{levinson1994introduction,boroditsky2011language}. Thinking for speaking, i.e. cognitive processes involved in the selection of words, their use in grammatical structures, and speech planning amongst others, has been shown to affect the habits of mind, impacting memory, attention, perception, and associations (\cite[22]{deutscher2010through}, \cite[72]{kramer2016feminist}). Such concepts are relevant for the present work, inasmuch as one of the differences in language on which scholars have concentrated is gender. This will be further explained in the next section.

%--------------------------------------
%	SUBSECTION 2
%-------------------------------------
\subsection{Gender in language} \label{sub:language/gender}
The relationship between language and gender is quite complex. It can be distinguished as linguistic (grammatical, referential, lexical) and extralinguistic gender (gender identity) \citep{corbett1991gender,unterbeck2011gender,cao-daume-iii-2020-toward}. Moreover, there is a distinction between linguistic and social gender \citep{unterbeck2011gender}. However, in the literature on language and gender, the latter term is used with different meanings (see for example \cite{nissen2002aspects}). Accordingly, this monograph opts for \textit{extralinguistic} gender that indicates the referents' gender identity. 

Grammatical gender is defined by \textcite[231]{hockett1958course} as ``classes of nouns reflected in the behavior of associated words''. It concerns the classification of nouns based on agreement between several words and the nouns. Word classes, such as articles, adjectives, and determiners, display inflection (see \cite{corbett1991gender}). Lexical gender refers to the semantic properties of words (lexical items) that are used to denote femaleness or maleness as in the case of nouns related to kinship, such as mother and father. Lexical gender is thus a property of the noun itself and not of its extralinguistic referent \citep{FUERTESOLIVERA2007219}. Referential gender relates to linguistic structures used to refer to the extralinguistic reality, i.e. to the gender identity of speakers and/or referents \citep{corbett1991gender,FUERTESOLIVERA2007219}. 

Linguistic gender finds different expression in languages, which can be classified into (i) grammatical gender; (ii) notional gender; and (iii) genderless languages (see Table \ref{tab:genderlanguage}) \citep{stahlberg2007representation}. Originally, the second category was named natural gender languages; however, due to the potential semantic overlap of natural with biological/anatomical sexual categories, the term notional is here preferred as in \textcite{mcconnell2013gender}. In grammatical gender languages, e.g. German and Italian, every noun has a gender, i.e. even inanimate objects and abstract concepts. Notional gender languages, e.g. Danish and English, are marked by lexical gender (girl/boy), which also influences compounds, e.g. chairman/chairwoman, derivational nouns, e.g. actor/actress, and gendered pronouns, e.g. she/he/it. In both (i) and (ii), gender is also referential. Human referents should thus be referred to with linguistic structures that are congruent with their gender identity. The relationship between linguistic gender, referential gender, and gender identity is not one-to-one (\cite{cao-daume-iii-2020-toward}). For instance, many European languages have binary structures, e.g. Italian \citep[99]{deutscher2010through} that hinder the representation of non-binary people. Accordingly, new linguistic forms must be introduced \citep{hornscheidt2021wie,itainclusivostoria}. %Derivational nouns might also carry gender inflections (e.g. actor/actress). 
Finally, in genderless languages, e.g. Turkish and Finnish, neither nouns nor pronouns express gender -- with a few exceptions such as references to kinship, e.g. \textit{sisko}, Finnish for sister. 



%------------------TABLE begins here --------------
\begin{table}[ht]
    \centering
    \caption[Expression of gender/sex in different language types]{Expression of gender/sex in different language types (\cite[166]{stahlberg2007representation})}
    \label{tab:genderlanguage}
\begin{tabularx}{\textwidth}{lQQQ}
        \lsptoprule
         &  {Grammatical  gender languages} &  {Notional gender  languages} &  {Genderless  languages} \\
         \midrule
         Frequency & High & Middle & Low \\
         Necessity & Often & Sometimes & Rare \\
         Linguistic Forms & Lexical, pronominal, grammatical & Lexical, pronominal & Lexical \\
         \lspbottomrule
    \end{tabularx}

\end{table}
%-----------------TABLE ends here --------------------

%In grammatical gender languages gender assignment for objects is formal (e.g. sun is masculine in Italian but feminine in German). For human referents assignment occurs on a semantic basis\footnote{This also occurs in notional gender languages in the case of lexical gender and gendered pronouns}. In linguistics, there is a widespread view that in such cases linguistic forms must be selected based on the referent's sex\footnote{Here it is argued that by the term sex the more appropriate expression sex assigned at birth is meant.} (e.g. \cite[34]{corbett1991gender},\cite{stahlberg2007representation}). Newer views recognize that in the case of people, assignment occurs based on sociocultural gender and on current ideas about sex and sexuality (\cite[9]{mcconnell2013gender}). Here it is important to underline that, even though terminology in this regard seems outdated and it is still discussed, gender identity is what counts in assignment of linguistic forms for human referents (\cite[153]{Zimman2019soclang}).

%The inclusion of non-binary people in discourse poses often a challenge due to the lack of appropriate or sufficient linguistic means. In European languages such as German\footnote{Although in German there is also the neuter gender, its use to refer to people is considered as degrading.}, Italian, and Spanish a binary distinction between female and male is dominant (\cite[99]{deutscher2010through}). Accordingly, new linguistic forms must be introduced (\cite{hornscheidt2021wie}, \cite{itainclusivostoria}). 

Additionally, gender in language (re)produces social expectations and norms. Although the representation of sex/gender varies across languages, \textcite[164]{stahlberg2007representation} note that maleness and femaleness can be always marked, thus concluding that sex/gender is perceived as an essential social category by speakers. The scholars also underline how such representations do not communicate a simple distinction between the sexes/genders. Rather, they reflect the underlying gender belief system, revealing evaluations and stereotypes in the way languages refer to the sex/genders. This is the reason why the question of gender in language goes well beyond grammar and is of great interest to linguists as well as to psychologists, sociologists and many others \citep[163]{corbett1991gender}. For instance, as further described in Subchapter \ref{Sexism in language}, female forms undergo \textit{semantic derogation} \citep[63]{lakoff1973language}, i.e. they take on an additional, negative meaning (e.g. mister/mistress). Gender-neutral forms, such as bigender nouns, can also trigger bias. In these cases, gender assignment usually occurs based on social gender, namely stereotypes by which a doctor is thought to be a man and a nurse a woman (see for example \cite{nissen2002aspects}). As already explained in Subsection \ref{sub:translatinggender}, differences among languages in gender structures and connotations are particularly challenging for translation and have therefore been one of the main focuses of feminist TS \citep{simon1996gender,von2016translation}.


Studies from psycholinguistics and cognitive linguistics on such a relationship between language and gender are reported in the next sections. Note that such research has been carried out quite exclusively focusing on European languages, i.e. in systems where the binary distinction between male and female is dominant \citep[199]{deutscher2010through}. Additionally, the differences among language structures and social connotations of gender are amongst the factors that contribute to issues of gender bias in MT along with, e.g., bias in the training datasets and in MT's functioning (algorithmic bias) (see Subsection \ref{sub:biasMT}).

Probably the first experiment to study the influence between language and gender was conducted at the Moscow Psychological Institute in 1915. Fifty participants were asked to visualise the days of the week as people and describe them. Participants described Monday, Tuesday, and Thursday as men but Wednesday, Friday, and Saturday as women. The only reason researchers could find for this is that the former are masculine and the latter are feminine in terms of grammatical gender in Russian \citep[209]{deutscher2010through}.

\textcite{konishi1993semantics} focused on two languages with grammatical gender, namely German and Spanish. To find whether grammatical gender carries connotations of femininity and masculinity, the researcher provided 40 German and 40 Mexican participants with a list of 54 words of opposed gender, such as bridge, i.e. \textit{die Brücke}, feminine in German, and \textit{el puente}, masculine in Spanish, in the two languages. Half of the words were masculine and half feminine. The participants were asked to give judgements using three semantic differential scales for evaluation, potency, and activity. The researcher found that masculine words received higher judgements on potency than female words in each language, concluding that grammatical gender influences perception and carries societal connotations of maleness and femaleness.

To understand whether participants indeed imagine objects as having stereotypically masculine or feminine qualities or whether the results are influenced by the naming of the object with the corresponding article, \textcite{phillips2003can} conducted a similar experiment. They asked German- and Spanish-speaking participants to describe objects that have opposite genders in the two languages. This time the study was carried out in English to test the influence of the first language even when not actively used. It resulted in participants using different adjectives according to the gender of the object. The word \textit{key} is masculine in German and feminine in Spanish. German speakers provided adjectives such as \textit{hard} or \textit{heavy}, whereas Spanish speakers preferred \textit{beautiful} or \textit{elegant}. This tendency was observed throughout the whole experiment, confirming that if objects are provided with gender, then people tend to mentally represent that object by highlighting its stereotypically masculine or feminine qualities \citep[65]{boroditsky2003sex}. %The starting point of such research was that some grammatical distinctions such as plural forms are easily observable in reality. When children learn a language presenting a grammatical gender system, they could actually believe that this grammatical property reflects significant differences between objects (\cite[929]{phillips2003can}). 

Another experiment by the scholars consisted of testing the memory of a group of Spanish and German speakers to see if gender exerts some sort of influence. Objects were given a masculine or feminine proper name. The gender of the name was consistent with the grammatical gender only half of the time. They found that participants tended to remember the object’s name when its grammatical gender and the gender of its proper name corresponded, e.g. Germans were more likely to remember Patrick than Patricia as a name for an apple since apple is masculine in German (\cite[929]{phillips2003can}, \cite[68]{boroditsky2003sex}). Results of this last experiment are considered particularly promising as, in this case, associations are not forced by asking participants to describe objects, thus bypassing the use of subjective judgements \citep[212--213]{deutscher2010through}.

\textcite{prewitt2012gendering} conducted an entirely different study in which they compared gender equality across 134 countries. Their purpose was to understand whether there is a correlation between gender equality and the typology of the language spoken, namely with grammatical gender, notional gender or genderless. Based on the results provided by \textcite{boroditsky2003sex}, the scholars wondered if the influence of gender was to be extended to the relationships between women and men as well \citep[269]{prewitt2012gendering}. The scholars compiled a list of 134 countries and compared the levels of gender equality by using the gender system of their official languages as metrics. They found that (i) higher levels of gender equality are found in countries where a notional gender language is spoken; (ii) countries where a grammatical gender language is spoken show lower levels of gender equality, and (iii) countries with genderless languages are to be found in the middle of this ranking. The authors hypothesise that the use of symmetrical gender forms, where men and women are equally visible, may promote gender equality. Grammatical gender languages often have asymmetrical gender forms, e.g. German where feminine nouns are derived from the masculine form by adding the suffix \textit{in}, and/or numerous gender suffixes, e.g. Italian where suffixes span from \textit{o}/\textit{a} to \textit{tore}/\textit{trice} and \textit{ore}/\textit{essa}. This complicates the creation and use of symmetrical gender-fair forms \citep[279]{prewitt2012gendering}. 

\textcite{santacreu2013female} created an index to measure how strong and common female/male distinctions are in languages using grammatical gender. Based on this, the scholars compiled a list of 84 countries to assess the effect of female/male distinctions on the political participation of women in such countries. Their findings were that (i) the more a language emphasises the aforementioned distinctions, the more the country in which it is spoken is likely to adopt a gender quota in politics; (ii) these very same countries also show a larger increase of women participating in the political agenda after the adoption of the quota. The explanation for (i) is that distinction in gender could trigger gender stereotypes, thus representing a barrier for women who would like to actively participate in politics. 

As explained in this section, languages can be classified as grammatical gender, notional gender or genderless based on their representation of sex/gender. Society and language seem to influence one another as regards gender marking: not only do speakers associate objects with stereotypically feminine or masculine characteristics according to the objects’ grammatical gender, but different gender systems can also be correlated with more or less gender equality. For instance, countries in which grammatical gender languages are spoken show lower levels of gender equality because, as research suggests, the constant distinction between masculine and feminine seems to foster gender stereotypes. This point is of particular interest. Through linguistic expressions and behaviour, gender identities can be (mis- or under-)represented to different extents in languages, hence reflecting and contributing to the prevailing sexism of Western societies. It emerges that language is a socially situated phenomenon that can contribute to discrimination. This aspect is further illustrated in the next section.

%--------------------------------------
%	SUBSECTION 3
%-------------------------------------
\subsection{Language as a social practice} \label{languapepowerdiscrimination}
In the previous sections, it was posited that linguistic differences can impact thought -- more specifically, memory, attention, perception, and associations. The focus was on gender in language, showing that gender differences in languages are often believed to reflect gender differences in society and vice versa (e.g. \cite{boroditsky2003sex}). Besides, gender marking can impact gender equality (e.g. \cite{santacreu2013female}). Language, as well as translation, is thus a form of social practice, both part of and conditioned by society \citep[55f]{fairclough2015language}. 

Consequently, it is not a neutral communication instrument, but rather a concrete action that shapes reality and reflects certain views of the world (\cite[7]{sprachhandeln2015tun}, \cite[79]{kramer2016feminist}). Language reflects social processes, power relations, and norms, but it can also create, normalise and reinforce them (\cite[55f]{fairclough2015language}, \cite[11]{sprachhandeln2015tun}, \cite[71]{kramer2016feminist}). This is also true for translation \citep{tymoczko2002introduction} as explained in Section \ref{TSrelevance}. People can be named, classified, categorised, evaluated and excluded through linguistic means. This effect is so profound that many of the existing naming possibilities seem evident and clear, and their construction, as well as their power, is not questioned \citep[11]{hornscheidt2008gender}.\footnote{Actually, \textcite[11]{hornscheidt2008gender} refers to German and Swedish in the passage, but such claims can be generalised.} Discrimination is hence (re)produced and normalised \citep[11f]{sprachhandeln2015tun}. 

This can be exemplified by analysing how assumptions on sex/gender are made in everyday life. A person’s gender identity and their identified pronouns are often assumed based on their voice, behaviour, clothing and so forth \citep[32]{prona2016mein}. This means that people tend to instinctively categorise a person as a woman or a man just by looking at them or by hearing their voice on the phone. In written communication, one's gender identity is also assumed by reading one's name. However, using such clues often leads to false assumptions and consequently to misgendering. Studies on misgendering show that it may cause feelings of stigma and psychological distress in transgender and non-binary people \citep{mclemore2015experiences,mclemore2018minority}. This demonstrates the importance of using linguistic forms that include and represent people of all genders in discourse. Moreover, reading gender through the physical aspect reinforces the gender binary, cisgenderism\footnote{The assumption that everyone is cisgender.} as well as norms and assumptions on how people and sexes/genders should be. Such norms and assumptions are reinforced by the use of binary gendered forms and also stereotypical expressions, e.g. \textit{both genders}, constantly implying that people identify either as women or as men. 

Gender is continuously produced through discourse in so many different ways that its production and conceptualisation are unquestioned. It is hence perceived and experienced as something natural \citep[24f]{hornscheidt2008gender}. This happens even before birth, when the doctor assigns sex after the inspection of the foetus' genitalia via ultrasound and pronounces sentences such as ``It's a girl''. Butler explains:
\begin{quote}
Consider the medical interpellation which (the recent emergence of the sonogram notwithstanding) shifts an infant from an ``it'' to a ``she'' or a ``he,'' 
%(sic) 
and in that naming, the girl is ``girled,'' %(sic) 
brought into the domain of language and kinship through the interpellation of gender. But that ``girling'' of the girl does not end there; on the contrary, that founding interpellation is reiterated by various authorities and throughout various intervals of time to reinforce or contest this naturalized effect. The naming is at once the setting of a boundary, and also the repeated inculcation of a norm. \citep[7--8]{butler2003bodies}
\end{quote}

\noindent This process by which sex is assigned through the pronunciation of the sentence ``it is a girl'' has far-reaching consequences for the subject in question. It is expected that they will identify as a woman and people around them will behave accordingly, for example by choosing a stereotypically feminine name, using feminine pronouns, buying stereotypically feminine clothes for them and so forth. Language reproduces, confirms and creates social norms \citep[11]{sprachhandeln2015tun}. 

Through generalisations and assumptions, linguistic discrimination can be explicitly or implicitly practised \citep[7]{sprachhandeln2015tun}. For this reason, when talking about people and sex/genders, one's language must be carefully analysed in terms of which sexes/genders are named, how people are categorised or evaluated based on their sexes/genders, and how they are referred to, i.e. with which linguistic form and/or expressions. Such linguistic decisions and actions are, more or less consciously, based on assumptions that (re)produce gender norms and expectations \citep[11]{sprachhandeln2015tun}. Through such analyses, language use can be questioned and modified to change the ruling norms, assumptions and expectations \citep[12]{sprachhandeln2015tun}. Language can ultimately become an instrument of resistance. In the same way, translators can decide whether to uncritically reproduce the sexism of source texts, openly highlight and criticise it, or remove it from target texts \citep{burton2010inverting}. %This is the main purpose of GFL as intended in this doctoral thesis: to question the sex/gender binary and foster the visibility of non-binary people.  

From a historical perspective, the interrelation between language and society has been a point of interest for feminist linguistics and TS since the 1960s (see for example \cite{cameron1985feminism,simon1996gender,kramer2016feminist,von2016translation}). During the second wave of feminism, feminists started focusing on language. Their starting point was that language is not only part of society, but that society is shaped by patriarchal notions. Language analysis should therefore help understand the system that oppresses women (see Subsection \ref{Sexism in language}) and their liberation might be achieved through language reform \citep[3]{cameron1985feminism}. Feminist language theory has focused on a variety of questions including: (i) the differences in language use between men and women; (ii) the representation of women in language through the use of stereotypes and metaphors; (iii) the reinforcement of gender norms and relationships through language \citep[8]{von2016translation}. Such work shows the strong link between linguistics and TS, focusing on language use as a social practice and representing the theoretical background of the present work. In the next section, studies on gender in the fields of cognitive and psycholinguistics are presented to provide concrete evidence for the impact of using gender-exclusive language.

%Views that language and society are intrinsic and influence one another are nowadays very widespread. Sociologist, linguists, philosophers as well as intellectuals propose different analyses on how a sexist use of language reflects and contributes to sexism in society (see \cite{priulla2014parole}, \cite{gasparrini2019non}, \cite{gheno2019potere}, \cite*{gheno2019femminili}, \cite{murgia2021stai}). Interestingly, this discourse is strongly developing in Italy, where, it is argued, the public debate on GFL has not been particularly intense (exceptions are the work of \cite{sabatini1986raccomandazioni} and \cite{robustelli2012linee}, \cite*{robustelli2014donne}). For instance, \textcite[89]{gasparrini2019non} claims that the way people express themselves reveals how they conceive the world, hence when sexist terms are used, people are being sexist. \textcite{priulla2014parole} analyses how the use of vulgar, sexist, homophobic (and often violent) words have become common in Italy, be it in television, in the social media or even in the political discourse. The sociologist states that tolerance towards such language led to people in the country being rather sexist and homophobic. Furthermore, \textcite{gheno2019femminili} summarises arguments that are commonly used against the feminisation of some high-level professions that are still very often used in the masculine form in Italian (e.g. engineer, minister, lawyer). She then quotes \textcite{sabatini1986raccomandazioni} who, in her guidelines against a sexist usage of the Italian language, affirmed that linguistic equality would foster societal equality. \textcite[15]{gheno2019femminili} argues indeed that through language we express our thoughts as well as our essence as human beings since language is what renders human beings as such (\cite[15]{gheno2019femminili}). Similar positions were adopted by \textcite{murgia2021stai} as well, who analyses sexist expressions used against women in everyday language in Italy, arguing that they help assess how unequal society is in the country. 

%What is also common in these analyses is the emerging urge for language reform and the adoption of GFL. In order to support such claim, in the next subchapter analyses on language sexism are reported followed by some studies in the field of cognitive linguistics. These show that the use of gender-exclusive language, specifically the use of male generics, have a negative impact on gender equality. The use of GFL is especially advocated by non-binary people. Non-binary activists also engage in language activism in order to be equally represented in language (\cite[34]{bergman2017non}). In English speaking countries, the debates include the use of gender-neutral title Mx by institutions and of singular \textit{they} by media outlets. 

To conclude, sexism is not the only form of discrimination that is constantly (re-)produced in language. Others include genderism, racism, ableism, and classism, amongst others \citep[11]{sprachhandeln2015tun}. These forms of discrimination cannot be regarded as distinct; they interact and overlap cumulatively. \textcite{crenshaw1989demarginalizing} called this phenomenon \textit{intersectionality}. An illustrative example of how intersectionality works is the lawsuit by Emma DeGraffenreid against General Motors (GM) in 1976. GM was accused of discriminating against black people but won the lawsuit. When the judges verified that the company hired black people and white women, they believed there was no discrimination against black women, ignoring that black women face discrimination precisely because they are both black and women \citep[142f]{crenshaw1989demarginalizing}. One analogy that can be used to understand intersectionality is that of a crossroads. Every road represents a line of identity, for example, gender, sexuality, race, and class. The unique combination (or intersection) of these roads describes the multiple oppression faced by a single person.


%--------------------------------------
%	SUBSECTION 4
%-------------------------------------
\subsection{Linguistic sexism} \label{Sexism in language}
%In the last subchapter, it was posited that societal beliefs and assumptions can be reflected in the way language is used. For instance, sexism and gender bias are (re-)produced by means of language (\cite[1]{menegatti2017gender}, \cite[5]{weatherall2002gender}). Language reflects gender stereotypes, thus reinforcing the belief system on which they are based and producing discrimination against people based, amongst others, on their gender/sex. As a consequence, language is biased and (re)produces discrimination in different ways that are discussed in this section.

As mentioned in the previous section, interest in sexism in language spread with second-wave feminism, i.e. by the mid/late 1960s. During that time, gender was introduced as an analytical category to account for the differences in societal roles and opportunities between men and women \citep[5]{von2016translation}. The feminist movement focused on the analysis of gender, which also manifests itself in language. Language as a powerful tool, used by institutions ruled by men, had to be reformed to free women from patriarchy and was therefore documented to reveal inequalities between men and women (\cite[73]{von1991feminist}, \cite[552f]{pauwels2003linguistic}, \cite[17]{von2016translation}).

As \textcite[167]{stahlberg2007representation} point out, a language proves to be sexist not because of the expression of gender, but rather because of how feminine forms are marked. In many languages the feminine form derives from the masculine one, hence being more complex \citep[167]{stahlberg2007representation}. An example is German, a grammatical gender language, where feminine nouns are often obtained by adding the suffix \textit{in} to masculine nouns, e.g. \textit{Nachbar} – \textit{Nachbar-in} to indicate the male and the female neighbour. In such cases, \textcite[167]{stahlberg2007representation} argue, men and women are not treated equally on a linguistic level. One could also suggest that derived feminine forms imply that women are somehow dependent on or subordinate to men. \textcite[8]{menegatti2017gender} cite other examples from Italian too. Although in Italian noun forms are often symmetrical with suffixes \textit{-o} and \textit{-a} referring to male and female, some others are not -- like the feminine suffix \textit{-essa} in \textit{dottoressa} as the feminine form of \textit{dottore} (doctor in English, a high-level profession). This phenomenon can also be found in English, as in the pairs \textit{hero-heroine}, and \textit{actor-actress}. 

Another main feature of a sexist use of language is the invisibilisation of women \citep[13]{weatherall2002gender}. Women are underrepresented in language due to the use of masculine generics. Masculine generics are linguistic forms used to refer to both men and to people in general \citep[169]{stahlberg2007representation}. This is especially common in idioms with grammatical gender like Italian or German, (e.g. \textit{studenti} or \textit{Studenten} is often used to refer to a mixed group of students) but is also common in notional gender languages like English with the use of the pronoun \textit{he} as generic \citep[9]{menegatti2017gender}. \textcite[175f]{stahlberg2007representation} summarise different studies to confirm that the use of masculine generics tends to evoke masculine associations, thus discriminating against women -- some of these can be found in Subsection \ref{sub:malegenerics}. The use of masculine generics has been an object of discussion in the last few years with different strategies proposed to avoid it. In German, generic masculine forms are nowadays less common in written language and usually gender-fair alternatives are chosen, e.g. \textit{Student*innen} or Student:innen (see \cite{boka2021richtlinie}). German-speaking countries have always been sensitive to gender issues and language policies have been implemented by institutions \citep[5]{sczesny2016can}. In English too, the generic \textit{he} is often avoided, and language authorities recommend the use of singular \textit{they} \citep{apasingularthey}. In Italian, the debate on GFL has always been less lively, focusing instead on the feminisation of agent nouns. For example, high-level professions such as \textit{minister} or \textit{president} are still often used in the masculine form even with reference to women (see \cite{gheno2019femminili}).


Another issue of sexism is the conceptual asymmetry in languages \citep[15]{litosseliti2014gender}. Several words describing women have no counterpart and vice versa. \textcite[5f]{menegatti2017gender} remark how this phenomenon relates to gender stereotypes according to which women are pure and family-oriented. Terms or expressions like \textit{working mother} or \textit{career woman} have no male counterpart. This asymmetry is also found in the field of high-level or male-dominated professions, i.e. job titles like policeman, fireman and so forth \citep[6]{menegatti2017gender}. As \textcite[553]{pauwels2003linguistic} underlines, the phenomenon of lexical gaps in this field is very telling. There is a lack of words describing women in certain roles and occupations whereas masculine lexical gaps in female-dominated professions are usually filled quickly.

Women are also often referred to in terms of their relationships to men \citep[65]{lakoff1973language}. They are characterised as someone’s wife or girlfriend. Besides, as \textcite[65f]{lakoff1973language} underlines, in journalistic writing men are more likely to be referred to just by their name, whereas women by relationship to a man. Forms of address even used to change depending on marital status. An unmarried woman was called \textit{Miss} while a married one  \textit{Mrs}. This distinction was and sometimes is still made in other languages such as German with \textit{Frau/Fräulein} and Italian \textit{signora/signorina}. This difference is crucial because it seems to imply that marital status was somehow more important for a woman than for a man because English only has one form for addressing a man, namely \textit{Mr.} \citep[21]{weatherall2002gender}. Since the 1960s, \textit{Ms} has been introduced to avoid such asymmetry with men, although the other forms of address are sometimes still used. Similarly, when telling stories about women in journal articles, reference is often made to their being mothers, even though this is irrelevant to the context of the story, to their womanhood \citep{murgia2021stai}, or, in general, to their appearance \citep[19]{weatherall2002gender}.

There are several further elements of a sexist use of language. These include metaphors. Metaphors used to describe women relate for example to immaturity, animals, clothing, food, vehicles, and furniture \citep{weatherall1999metaphorical}. This means that women are often referred to as \textit{babe}, \textit{bitch}, \textit{sweetie pie} and so forth. In particular, the animals named in such metaphors are usually domesticated or hunted for sport \citep{weatherall1999metaphorical}. Other forms of linguistic sexism are the lack of expressions to refer to the experience of women, especially in the field of sex and sexuality, giving compliments by comparing them to men, e.g. ``she thinks like a man'', insulting men by comparing them to women, e.g. ``he's a real old woman'' (see \cite{weatherall2002gender}), and the use of gender stereotypes (see \cite{apagenderenglish}).


%--------------------------------------
%	SUBSECTION 5
%-------------------------------------
%https://onlinelibrary.wiley.com/action/authenticateSharedSP
%Jeanette Silveira. 1980. Generic Masculine Words and Thinking. --> no access 
\subsection{The societal impact of gender-exclusive language} \label{sub:malegenerics}
Feminist analyses of sexist language use show how women are linguistically disregarded in a multitude of ways. Many claim that this also has an impact on society (see for example \cite{gasparrini2019non} and \cite{murgia2021stai}). In particular, research in psycholinguistics and/or cognitive linguistics has concentrated on the use of masculine generics and its impact on (i) the perception that the prototypical human is a man; (ii) on the presence of women in the labour market; and (iii) their visibility. Some studies in this field are reported in this section. This overview is by no means complete, but it demonstrates the concrete societal impact of gender-exclusive language. Accordingly, it is time to adopt GFL strategies to make all (or no) genders visible and hence foster linguistic and social equality. 

Based on previous findings, \textcite{hamilton1988using} wondered whether the use of masculine generics also had an impact on speakers rather than only on listeners. The scholar asked 120 female and male college students to complete a set of sentences in two different conditions, namely, using masculine generics or gender-fair alternatives. The participants then had to describe the images evoked in their minds. Bias was present in both conditions but to a far greater extent in the first one, i.e. when completing sentences with masculine generics. Moreover, men were generally found to be more biased than women.

\textcite{bem1973does} discovered that women hesitated to apply for jobs if the job adverts were presented solely in masculine forms. They carried out an experiment asking female and male students to indicate how likely they were to apply for twelve stereotypically male or female jobs. The advertisements were written (i) in gender-exclusive language; (ii) in (binary) gender-inclusive language; (iii) by exclusively appealing to the gender that rarely takes such jobs. Gender-exclusive language affected both female and male students, who were not interested in jobs if the advertisement targeted the opposite gender. In addition, women were more likely to apply for a job if the advert addressed them explicitly. 

\textcite{stahlberg2001effekte} focused on the effect of masculine generics on German native speakers in different studies. Participants (i) were asked to fill out a questionnaire about personal preferences, in which they had to name their favourite athlete, musician and so forth. Questions were phrased using either only masculine forms, paired forms, or gender-neutral forms. Similarly, (ii) other participants were asked about their media consumption and, among the questions, had to name three athletes, singers, politicians, etc. they knew. In this case, the questionnaire was written using masculine forms, paired forms or the nowadays outdated \textit{Binnen-I}. \footnote{The so-called \textit{Binnen-I}, i.e. a capital I separating the masculine from the feminine form of the nouns, is nowadays considered as exclusive as it implicitly reproduces a binary conception of gender (see for example \cite[49]{hornscheidt2021wie}).} Furthermore, (iii) \textcite{stahlberg2001effekte} asked participants to name possible candidates for the office of chancellor in Germany, presenting questions either using the masculine form or paired forms. In every study configuration, women were named less often when masculine forms only were used.

\textcite{stout2011he} conducted three studies assessing the reaction of women and men to the use of gendered language in a professional context. Participants were asked to read (study \#1) a job description or participate in a mock interview where they were given a description of the role and the work environment (study \#2 and \#3). The descriptions were delivered using (i) only masculine forms (all studies), (ii) paired forms (study \#1), and (iii) also neutral forms (study \#2 and \#3). Participants were asked about perceived sexism, feelings of inclusion or exclusion in the work environment, motivation to pursue the job, and identification with it. The researchers found that, although both men and women perceived the jobs described in gender-exclusive terms as more sexist, women felt generally more ostracised, less motivated and identified less with such jobs than men. In a nutshell, women were more affected by the use of gender-exclusive language than men.

\textcite{horvath2016reducing} carried out a study to investigate the effect of masculine generics on the evaluations of women in hiring simulations. They asked participants to evaluate male and female applicants for two different job positions, namely a low-status role and a high-status leadership position. The advertisements for the positions were written using masculine forms, masculine forms with the addition of ``m/f'' in brackets, or paired forms. \textcite{horvath2016reducing} found that women were perceived as similarly suitable for the lower status position in each case, but were perceived as less suitable for the higher status position when masculine forms, with or without ``m/f'' in brackets, were used.

Noting that white men are often leaders in different domains, \textcite{bailey2017counts} wanted to find whether this is supported by an androcentric belief system. The researchers conducted an experiment and asked participants to first pre-select three faces out of a set of 20 photos and then choose a single face that could represent humanity as a whole. The task was phrased using either a gender-exclusive, a gender-neutral, or a binary expression, namely (i) mankind, (ii) human, or (iii) man or woman. \citeauthor{bailey2017counts} observed an androcentric bias in each case, i.e. participants generally selected a male person as representative of humanity. However, language had an impact on the selection, since in (iii), participants' choices were more balanced.

These studies represent only a very limited selection of research on masculine generics in the field of cognitive linguistics and psycholinguistics. The consequences of gender-exclusive language are nevertheless quite clear. Masculine generics foster the idea that the prototypical human is a man, who is often a straight, white, middle-class, able-bodied cis-man, hence men being the norm and women being rather the exception. Women are less visible, with even `real-life' consequences such as their exclusion from specific job positions. Non-binary individuals and non-binary GFL are still underresearched. Some first studies are presented in Chapter \ref{sec:impactgfl}.

This chapter has elucidated the close connection between TS and sociolinguistics. Language does not exist in isolation; it is shaped by societal assumptions and norms, which influence language use. These norms and assumptions are also reflected in translation practices, as demonstrated by feminist and queer translation scholars. Consequently, translators bear social responsibility for ensuring respectful representation of marginalised groups, such as non-binary individuals. This social responsibility of translators underscores that translation extends beyond a mere process of linguistic transfer, highlighting its potential to foster societal change.
